% from ega0-1-fr.tex
\oldpage[0\textsubscript{I}]{37}
\[
\begin{tikzcd}[column sep=large,row sep=large]
\mathcal{C}\otimes_{\mathcal{A}}\mathcal{C}\otimes_{\mathcal{A}}\mathcal{C}
  \arrow[r,"\varphi\otimes 1"]
  \arrow[d,"1\otimes\varphi"']
&
\mathcal{C}\otimes_{\mathcal{A}}\mathcal{C}
  \arrow[d,"\varphi"]
\\
\mathcal{C}\otimes_{\mathcal{A}}\mathcal{C}
  \arrow[r,"\varphi"']
&
\mathcal{C}
\end{tikzcd}
\]
soit commutatif~; 2\textsuperscript{o} pour tout ouvert $U\subset X$ et toute section
$s\in\Gamma(U,\mathcal{C})$, on ait
$\varphi((e|U)\otimes s)=\varphi(s\otimes(e|U))=s$. Dire que
$\mathcal{C}$ est une $\mathcal{A}$-Algèbre commutative revient en outre à
dire que le diagramme
\[
\begin{tikzcd}[column sep=large,row sep=large]
\mathcal{C}\otimes_{\mathcal{A}}\mathcal{C}
  \arrow[rr,"\sigma"]
  \arrow[dr,"\varphi"']
&&
\mathcal{C}\otimes_{\mathcal{A}}\mathcal{C}
  \arrow[dl,"\varphi"]
\\
& \mathcal{C} &
\end{tikzcd}
\]
est commutatif, $\sigma$ désignant la symétrie canonique du produit tensoriel
$\mathcal{C}\otimes_{\mathcal{A}}\mathcal{C}$.

Les homomorphismes de $\mathcal{A}$-Algèbres se définissent aussi comme les
homomorphismes de $\mathcal{A}$-Modules dans (G, II, 2.2), mais naturellement
ne forment plus un groupe abélien.

Si $\mathcal{M}$ est un sous-$\mathcal{A}$-Module d'une
$\mathcal{A}$-Algèbre $\mathcal{C}$, la sous-$\mathcal{A}$-Algèbre de
$\mathcal{C}$ engendrée par $\mathcal{M}$ est la somme des images des
homomorphismes $\bigotimes^n\mathcal{M}\to\mathcal{C}$ (pour les
$n\geqslant 0$). C'est aussi le faisceau associé au préfaisceau
$U\mapsto\mathcal{B}(U)$ d'algèbres, $\mathcal{B}(U)$ étant la sous-algèbre
de $\Gamma(U,\mathcal{C})$ engendrée par le sous-module
$\Gamma(U,\mathcal{M})$.
\end{env}

\begin{env}[4.1.4]
\label{0.4.1.4-fr}
On dit qu'un faisceau d'anneaux $\mathcal{A}$ sur un espace topologique $X$
est \emph{réduit en un point} $x$ de $X$ si la fibre $\mathcal{A}_x$ est un
anneau \emph{réduit} (1.1.1)~; on dit que $\mathcal{A}$ est \emph{réduit}
s'il est réduit en tout point de $X$. Rappelons qu'un anneau $A$ est dit
\emph{régulier} si chacun des anneaux locaux $A_{\mathfrak p}$ (où
$\mathfrak p$ parcourt l'ensemble des idéaux premiers de $A$) est un anneau
local régulier~; nous dirons qu'un faisceau d'anneaux $\mathcal{A}$ sur $X$
est \emph{régulier en un point} $x$ (resp. \emph{régulier}) si la fibre
$\mathcal{A}_x$ est un anneau régulier (resp. si $\mathcal{A}$ est régulier
en tout point). Enfin, nous dirons qu'un faisceau d'anneaux $\mathcal{A}$ sur
$X$ est \emph{normal en un point} $x$ (resp. \emph{normal}) si sa fibre
$\mathcal{A}_x$ est un anneau intègre et intégralement clos (resp. si
$\mathcal{A}$ est normal en tout point). Nous dirons qu'un espace annelé
$(X,\mathcal{A})$ a l'une des propriétés précédentes si le faisceau d'anneaux
$\mathcal{A}$ a cette propriété.

Un faisceau d'anneaux \emph{gradués} $\mathcal{A}$ est par définition un
faisceau d'anneaux qui est somme directe (G, II, 2.7) d'une famille
$(\mathcal{A}_n)_{n\in\mathbf Z}$ de faisceaux de groupes abéliens avec les
conditions $\mathcal{A}_m\mathcal{A}_n\subset\mathcal{A}_{m+n}$~; un
$\mathcal{A}$-\emph{Module gradué} est un $\mathcal{A}$-Module
$\mathcal{F}$ somme directe d'une famille
$(\mathcal{F}_n)_{n\in\mathbf Z}$ de faisceaux de groupes abéliens,
satisfaisant aux conditions
$\mathcal{A}_m\mathcal{F}_n\subset\mathcal{F}_{m+n}$. Il revient d'ailleurs
au même de dire que l'on a
$(\mathcal{A}_m)_x(\mathcal{A}_n)_x\subset(\mathcal{A}_{m+n})_x$
(resp.
$(\mathcal{A}_m)_x(\mathcal{F}_n)_x\subset(\mathcal{F}_{m+n})_x$)
en tout point $x$.
\end{env}

\begin{env}[4.1.5]
\label{0.4.1.5-fr}
Étant donné un espace annelé $(X,\mathcal{A})$ (non nécessairement
commutatif), nous ne rappellerons pas ici les définitions des bifoncteurs
$\mathcal{F}\otimes_{\mathcal{A}}\mathcal{G}$,
$\mathcal{H}\!om_{\mathcal{A}}(\mathcal{F},\mathcal{G})$ et
$\operatorname{Hom}_{\mathcal{A}}(\mathcal{F},\mathcal{G})$
(G, II, 2.8 et 2.2) dans les catégories des $\mathcal{A}$-Modules à gauche ou
à droite (selon les cas), à valeurs dans la catégorie des faisceaux de groupes
abéliens (ou plus

